\subsection{A weaker criterion from the deck action on two-primary torsion}

\begin{claim}{A.C2.7}{Proved}
Let $C$ be a smooth proper geometrically connected curve of genus at
least two over a field $K$ finitely generated over $\F_2$. Suppose that,
over $\overline K$, it admits a connected nontrivial finite \'etale
Galois cover $D\to C$ of odd degree, with group $G$, for which
\[
             G\longrightarrow
             \operatorname{Aut}_{\F_2}(J_D[2](\overline K))
                                                               \tag{C2.19}
\]
is not faithful. Then the geometrically prime-to-two point-section
map of $C$ is injective. The conclusion also holds for its
finite-extension forms and hyperbolic open subcurves.

More precisely, a pair of distinct points giving equal section classes
would force (C2.19) to be faithful for every such cover. If
$m_2(G)$ denotes the least dimension of a faithful representation of
$G$ over $\F_2$, the same hypothetical pair would force
\[
                            f(D)-f(C)\ge m_2(G).               \tag{C2.20}
\]
In particular, a single such cover with $f(D)=f(C)$ suffices for
point-section injectivity.
\end{claim}
\begin{proof}
Work first over an algebraic closure. The curve $D$ is not
hyperelliptic. Otherwise its hyperelliptic involution is unique,
so every deck transformation descends to its genus-zero quotient.
Choose an element $\sigma\in G$ of odd prime order. Its image on
$\mathbf P^1$ has the same order, since the kernel of descent is
the order-two hyperelliptic group. A finite odd-order automorphism
of $\mathbf P^1$ in characteristic two has a fixed geometric
point. Its fibre on $D$ has at most two geometric points, so
$\sigma$ fixes a point in that fibre. This contradicts the free
deck action of an \'etale Galois cover.

Let $P',Q'\in D(\overline K)$ lie over distinct points of $C$.
The class $T=[Q'-P']$ has trivial stabilizer in $G$. Indeed,
$\delta T=T$ for $1\ne\delta\in G$ would give
\[
                     \delta Q'+P'\sim Q'+\delta P'.
                                                               \tag{C2.21}
\]
The two effective degree-two divisors in (C2.21) are distinct:
their terms lying over each of the two base points would
otherwise force both $\delta P'=P'$ and $\delta Q'=Q'$.
A linear equivalence between distinct effective degree-two
divisors on a curve of genus at least two gives a degree-two
map to $\mathbf P^1$, after removing any common base point.
A degree-one map is excluded by the genus. Thus (C2.21) would
make $D$ hyperelliptic, a contradiction.

The \'etale two-adic Tate module of $J_D$ is free of rank $f(D)$,
and its geometric two-primary torsion is the quotient of its
rational Tate module by this lattice. An odd-order automorphism
acting trivially on $J_D[2](\overline K)$ acts trivially on the
whole two-adic module. This can be checked without any
torsion-free congruence-group citation: if a matrix $A$ of odd
order $m$ satisfies $A\equiv1\pmod{2^r}$, then
\[
 A^m\equiv1+m(A-1)\pmod{2^{r+1}}
\]
implies $A\equiv1\pmod{2^{r+1}}$. Starting at $r=1$ and
iterating proves $A=1$. Thus any element in the kernel of
(C2.19) fixes every two-primary torsion point of $J_D$.

Now suppose that distinct $P,Q\in C(K)$ give equal section
classes. Descend the cover and its deck transformations after
a finite extension of $K$, and extend again to make a chosen
lift $P'$ rational. By A.3 there is a lift $Q'$ with
$T=[Q'-P']$ two-primary. A nonidentity element in the kernel
of (C2.19) would fix $T$, contrary to its trivial stabilizer.
This proves the criterion and the necessity of faithfulness.
The same argument after restriction proves the stated form and
open-subcurve variants.

For the quantitative assertion, the averaging idempotent
$|G|^{-1}\sum_{g\in G}g$ acts integrally on the two-adic Tate
module, since $|G|$ is odd. Its invariant summand is the
pullback of the two-adic Tate module of $J_C$: pullback
followed by norm is multiplication by $|G|$, and norm followed
by pullback is the sum of deck transformations. This summand
has rank $f(C)$. The complementary summand has rank
$f(D)-f(C)$, and its reduction modulo two must be a faithful
$G$-representation because the invariant summand is trivial.
This proves (C2.20).
\end{proof}

\begin{claim}{A.C2.8}{Proved}
In the genus-two cyclic situation of A.C2.5, let $S$ be the
union of the nonzero doubling orbits that miss $\Theta$.
The kernel of the deck action on geometric two-primary torsion
has order
\[
                    d=\gcd\bigl(n,\{i:i\in S\}\bigr),        \tag{C2.22}
\]
where $d=n$ if $S$ is empty. Consequently $d>1$ already implies
point-section injectivity by A.C2.7.

In particular, if $n=\ell^a$ is an odd prime power, it suffices
that every doubling orbit on $(\mathbf Z/n)^\times$ meets
$\{i:L^i\in\Theta\}$. No condition is needed on characters
of smaller order. The simple sufficient condition is
\[
 L\in\Theta\text{ of exact order }\ell^a,\qquad
                 \langle2,-1\rangle=(\mathbf Z/\ell^a)^\times.
                                                               \tag{C2.23}
\]
For a cyclic group of odd prime-power order $\ell^a$ one has
$m_2(G)=\operatorname{ord}_{\ell^a}(2)$.
\end{claim}
\begin{proof}
The Artin--Schreier description of \'etale $H^1$ identifies its
scalar extension with the stable Frobenius part of
$H^1(D,\mathcal O_D)$. The calculation of A.C2.5 therefore
shows that its nontrivial cyclic characters are exactly those
indexed by $S$, with multiplicity one. Dualizing to geometric
two-torsion may invert all characters; it does not change
their common kernel. A power $\delta^k$ acts trivially exactly
when $n\mid ik$ for every $i\in S$. The number of such powers
is the greatest common divisor in (C2.22).

If $n=\ell^a$, this kernel is nontrivial exactly when no
primitive character survives, since a primitive index is
coprime to $n$, whereas all nonprimitive indices are divisible
by $\ell$. Symmetry of $\Theta$ gives (C2.23) as in A.C2.5.
For example at $n=9$, a theta point of exact order nine kills
the single orbit of all six primitive characters. The
order-three block may remain ordinary, but the deck action
then has a kernel of order three, which suffices here. This
does not assert existence of such a nine-torsion theta point
on every curve.

Finally, an irreducible $\F_2$-representation of a cyclic
group corresponds to a doubling orbit of its geometric
characters. A faithful character of order $\ell^a$ has orbit
length $\operatorname{ord}_{\ell^a}(2)$. Any faithful
representation of this cyclic prime-power group must contain
such a character: otherwise all character orders divide
$\ell^{a-1}$ and a subgroup of order $\ell$ acts trivially.
One faithful orbit already gives a faithful representation,
proving the formula for $m_2(G)$.
\end{proof}

The full-quotient obstruction in A.C2.6 and the weaker deck
criterion are compatible. Surviving order-three characters
prevent the entire Prym from having $2$-rank zero, but at
degree $3^a$, $a>1$, they need not make the deck action
faithful. Torsion density still does not guarantee the
prime-power order or simultaneous primitive-orbit
intersections required by (C2.23). A degree-sized lower
bound in place of (C2.20), or a construction forcing a
sufficiently small normalized $2$-rank in general towers,
has not been established here.

\subsection{What the entire first cubic tower still cannot detect}

\begin{claim}{A.C2.9}{Proved; an obstruction to the specified target method}
Let $C/\F_4$ be the ordinary genus-two curve
\[
 y^2+y=x+1/x+\alpha/(x+1),\qquad \alpha^2+\alpha+1=0.
\]
Its full mod-three abelian \'etale cover $U\to C$ has group
$(\Z/3)^4$, degree $81$, genus $82$, and ordinary Jacobian.
Every intermediate curve is ordinary. Thus neither $U$ nor an
intermediate curve admits a nonconstant map to an abelian variety
of two-rank zero. The deck action on geometric two-primary torsion
is faithful for every nontrivial intermediate Galois cover of $C$.
The quotient of $U$ by a lifted hyperelliptic involution has genus
$40$ and is ordinary. These statements concern the complete first
cubic compositum; they do not imply the corresponding assertions
for higher three-power covers and do not refute A1.
\end{claim}
\begin{proof}
Work geometrically, and base the Abel map at a Weierstrass point.
Its pullback of $[3]:J_C\to J_C$ is connected because the Abel
map induces the maximal abelian quotient of the geometric
fundamental group \citep[III, Proposition 9.1]{A_Milne2008}.
Hence $G=\Gal(U/C)=(\Z/3)^4$ and Riemann--Hurwitz gives
$g(U)=1+81=82$. There are $40$ index-three subgroups $K$ of
$G$, and $D_K=U/K$ runs through all the connected cyclic cubic
covers. Write $P_K$ for the connected Prym of $D_K\to C$.
Claim A.C2.2 proves
\[
  J_{D_K}\sim J_C\times P_K,\qquad P_K\sim E_K^2,
\]
where $E_K$ is an ordinary elliptic reflection quotient and
$\dim P_K=2$.

Here is an exact exhaustion of the Jacobian factors. For a subgroup
$H\subset G$ put $e_H=|H|^{-1}\sum_{h\in H}h$ in $\Q[G]$.
The idempotents $e_G$ and $e_K-e_G$, for the $40$ hyperplanes
$K$, are pairwise orthogonal. For distinct $K,L$, one has
$e_Ke_L=e_G$ because $KL=G$. Every nonzero vector of $\F_3^4$
lies in $13$ hyperplanes, whereas zero lies in $40$, so
\[
 \sum_K\sum_{h\in K}h=27\,1+13\sum_{g\in G}g,
 \qquad e_G+\sum_K(e_K-e_G)=1.
\]
Pullback and norm identify these factors with $J_C$ and the
pullbacks of the $P_K$. Their sum therefore gives an isogeny
\[
                  J_U\sim J_C\times\prod_{40}E_K^2.       \tag{C2.24}
\]
All denominators in the displayed idempotents are powers of three,
so they split the two-adic Tate module integrally as well; no
two-primary component is omitted. The factor dimensions total
$2+40\cdot2=82$.

Ordinariness is preserved under products and isogenies, so $U$ is
ordinary. Pullback and norm make the Jacobian of every intermediate
curve an isogeny factor of $J_U$, hence ordinary. The image of a
homomorphism from an ordinary abelian variety to one of two-rank
zero has dimension zero: it is both an ordinary quotient and a
two-rank-zero subvariety. The Albanese property proves the assertion
about maps of curves.

Each nontrivial character pair of $G$ occurs in exactly one $P_K$.
Its kernel is $K$. Since $\bigcap_K K=1$, these characters act
faithfully on geometric two-torsion. For an intermediate cover
$U/H\to C$, every nontrivial character of $G/H$ still occurs,
with the same conclusion. Thus the sufficient criterion A.C2.7
does not apply to these covers.

The involution lifts on the Abel pullback as
$(P,z)\mapsto(\iota P,-z)$, and inverts $G$. On the fibre over
each of the three hyperelliptic branch points it is an affine
inversion of the odd $G$-torsor, with exactly one fixed point.
The local \'etale map to $C$ preserves its different exponent two.
The total different is six, whence
\[
 162=2\bigl(2g(U/\langle\tau\rangle)-2\bigr)+6
\]
and the quotient genus is $40$. Its Jacobian is again an ordinary
isogeny factor of $J_U$.

The finite integer certificate
\path{code/A_char2_cubic_tower_certificate.py} verifies all $40$
hyperplanes, their incidence counts, $861$ group-ring convolution
identities, the character partition, and the unique fixed point of
every affine inversion. The proof above supplies the identification
of these group-ring calculations with actual Jacobian factors.
\end{proof}

\begin{claim}{A.C2.10}{Proved; the inertia obstruction in equivariant odd towers}
Let $C$ be any genus-two curve over an algebraically closed field
of characteristic two. Let $V\to C$ be an odd-degree \'etale
Galois cover equivariant under the hyperelliptic involution, so
that $V\to\mathbf P^1$ has Galois group $\Gamma$ in an exact
sequence
\[
             1\longrightarrow N\longrightarrow\Gamma
               \longrightarrow C_2\longrightarrow1.
\]
Every map between intermediate curves of this $\Gamma$-cover
has ramification indices one or two. No odd-order subgroup of
$\Gamma$ has a genus-one quotient of $V$. Consequently the
usual elliptic group quotients in this tower cannot supply the
odd common ramification index in A.21 or a prime-to-two Galois
elliptic quotient for A.29l. This statement does not exclude maps
or additional automorphisms outside the specified tower group.
\end{claim}
\begin{proof}
Since $V\to C$ is \'etale, the ramification indices of
$V\to\mathbf P^1$ are those of $C\to\mathbf P^1$, namely
one and two. If $I_R$ is inertia at a point of $V$, the inertia
of $V\to V/H$ is $I_R\cap H$ for any subgroup $H\subset\Gamma$.
All intermediate ramification indices therefore divide two.

An odd-order subgroup $H$ maps trivially to $C_2$, so lies in
$N$. The cover $V/H\to C$ is \'etale, of degree $[N:H]$, and
\[
                 g(V/H)=1+[N:H](g(C)-1)\ge2.
\]
This proves both obstructions. Equivariant refinements are available
for any given odd Galois cover: intersect its normal subgroup of
the geometric fundamental group with the transform under the
hyperelliptic involution. The resulting odd quotient is invariant.
One can realize a lifted involution by normalizing the fibre product
with the conjugate cover and taking a component fixed by factor
interchange; the number of components is odd. Thus the obstruction
persists after these refinements, without asserting anything about
arbitrary maps outside their group quotients.
\end{proof}

\subsection{An exact nine-torsion point resolves the cubic-tower example}

\begin{claim}{A.C2.11}{Proved}
The curve
\[
 C/\F_4:\quad y^2+y=x+1/x+\alpha/(x+1),
 \qquad\alpha^2+\alpha+1=0,
\]
has an injective geometrically prime-to-two point-section map over
every finite extension. The conclusion also holds for its
finite-extension forms over fields finitely generated over $\F_2$
and for its hyperbolic open subcurves. An explicit cyclic \'etale
degree-nine cover has genus ten, two-rank four, and an order-three
deck subgroup acting trivially on all geometric two-primary torsion.
This higher cover overcomes the valid first-cubic obstruction
A.C2.9; the entire Prym is not required to have two-rank zero.
\end{claim}
\begin{proof}
Let $W_0,W_1,W_\infty$ be the unique points above $0,1,\infty$.
The different exponent at each simple Artin--Schreier pole is two,
so
\[
 \operatorname{div}(dx)=2W_0+2W_1-2W_\infty,
 \qquad B=W_0+W_1-W_\infty.
\]
The line bundle of locally exact differentials is $\mathcal O(B)$:
locally its transition functions square to those of ordinary
differentials, and the exact differential $dx$ has the displayed
divisor. The Raynaud theta is consequently
\[
                 \Theta=\{\mathcal O(R-B):R\in C(\overline\F_2)\}.
\]

Choose $t$ with $t^4+t^3+t^2+t+1=0$ and put
$\alpha=t+t^{-1}=t^3+t^2+1$. Then $\alpha^2+\alpha+1=0$.
Choose $\beta$ with $\beta^2+\beta=t+1$.
The identity
\[
                 t+t^{-1}+\alpha/(t+1)=t+1
\]
makes $R=(t,\beta)$ a point of $C$. These choices lie in
$\F_{16}$ and $\F_{256}$ respectively; the absolute trace of
$t+1$ from $\F_{16}$ is one.

Define
\[
\begin{aligned}
 b(x)&=x^3+(t^2+t+1)x^2+(t^2+t+1)x+t^2,\\
 a(x)&=x^3+t x^2+t x+t^2,\\
 h&=a(x)+b(x)(y+\beta),\qquad u=x(x+1).
\end{aligned}
\]
Exact polynomial multiplication modulo $t^4+t^3+t^2+t+1$ gives
\[
 \begin{aligned}
 u(a^2+ab)+b^2\bigl(x^3+x^2+(1+\alpha)x+1+(t+1)u\bigr)
                          &=(x+t)^9,\\
 \operatorname{Nm}_{C/\mathbf P^1}(h)&=(x+t)^9/u.
 \end{aligned}                                                   \tag{C2.25}
\]
Here $a(t)=0$ and $b(t)=1$, so $h(R)=0$ and its conjugate
value at $R$ is one. Thus all nine zeros lie at $R$.
The poles are exactly simple at $W_0,W_1$ and of order seven
at $W_\infty$: $b(0)=t^2$ and $b(1)=1+t^2$ are nonzero,
and $x^3y$ has the unique largest pole order seven at infinity.
Therefore
\[
 \operatorname{div}(h)=9R-W_0-W_1-7W_\infty,
 \qquad \operatorname{div}(h/u^4)=9(R-B).                         \tag{C2.26}
\]
The second equality uses
$\operatorname{div}(u)=2W_0+2W_1-4W_\infty$.

Hence $\eta=\mathcal O(R-B)$ lies on $\Theta$ and is killed
by nine. It is not killed by three. Claim A.C2.2 and the exact
theta formula A.C2.5 give $\Theta\cap J_C[3]=\varnothing$:
every cubic Prym is ordinary since
$(1+1+\alpha)^3+\alpha=1+\alpha\ne0$, and zero is absent
from $\Theta$ because $C$ is ordinary. Thus $\eta$ has exact
order nine.

Let $D\to C$ be the connected cyclic \'etale degree-nine cover
associated with $\eta$, obtained by taking a ninth root of $h/u^4$.
Its genus is ten. The primitive characters form the single doubling
orbit $1,2,4,8,7,5$ modulo nine. Since $\eta\in\Theta$, Frobenius
is nilpotent on their six one-dimensional cohomology spaces.
The orbit $3,6$ survives, since there is no nonzero three-torsion
on $\Theta$. Formula (C2.16) gives $f(D)=2+2=4$.
The common kernel of the surviving nontrivial characters $3,6$
is the order-three subgroup of the deck group. By A.C2.8 it
acts trivially on every geometric two-primary torsion point.

A section collision on $C$ lifts by A.3 to a two-primary difference
$[Q'-P']$ on $D$. A nonidentity element $\sigma$ of this
order-three kernel fixes that difference, giving
\[
                     Q'+\sigma P'\sim P'+\sigma Q'.
\]
The effective degree-two divisors are distinct, since the lifts
lie over distinct base points and the deck action is free. This
would make $D$ hyperelliptic, contrary to A.C2.7. This proves A1 for
the stated curve.
All data descend after a finite extension; the same lifting and
Kummer argument proves the finitely generated and open-subcurve
variants.

The exact certificate \path{code/A_theta_nine_certificate.py}
checks (C2.25), every local value, the trace and the two character
orbits. Root independently verified the polynomial identity with a
separate quotient-polynomial implementation in
\path{code/A_theta_nine_independent.py}; its full mathematical
review is \path{audit/A_theta_nine_root_review.md}.
\end{proof}

\subsection{Prime-power theta torsion is confined to a finite division field}

\begin{claim}{A.C2.12}{Proved; a finite search bound, not an existence theorem}
Let $J/\F_q$ be a genus-two Jacobian and let $\Theta$ be a
geometrically irreducible theta divisor defined over $\F_q$.
Let $\ell$ be an odd prime different from the characteristic,
and choose $r$ such that $J[\ell]$ is rational over $\F_{q^r}$.
Then
\[
 \Theta(\overline\F_q)\cap J[\ell^\infty]
       \subset J(\F_{q^r})[\ell^\infty].                         \tag{C2.27}
\]
In particular the intersection is finite. For the Raynaud theta
of the curve in A.C2.11, every three-primary point is killed by
$27$. Since theta three-torsion is absent, its only possible
nonzero three-primary orders are nine and $27$; existence of
order nine was proved separately, and order $27$ is not inferred.
\end{claim}
\begin{proof}
This is a direct genus-two form of the bounded-prime-support
phenomenon for almost rational torsion
\citep[Proposition 2, p.16]{A_BoxallGrant2006}. We give the
fixed-field proof, including the case $\ell=3$.

Write $\pi$ for $q$-power Frobenius and $\sigma=\pi^r$.
On $T_\ell J$, $\sigma$ is congruent to one modulo $\ell$.
Suppose $\eta\in\Theta\cap J[\ell^\infty]$ is not fixed.
The principal congruence subgroup modulo $\ell^a$ is an
$\ell$-group, so there is a least $b\ge1$ with
$\sigma^{\ell^b}\eta=\eta$. Put
$\tau=\sigma^{\ell^{b-1}}=1+\ell M$ on $T_\ell J$.
The operator
\[
 U=\frac{1+\tau+\cdots+\tau^{\ell-1}}\ell
   =1+\sum_{j=1}^{\ell-1}
       \binom{\ell}{j+1}\ell^{j-1}M^j
\]
is integral and congruent to one modulo $\ell$. At $\ell=3$
the final coefficient is three, so this case is included.
Thus $U$ acts invertibly on geometric $\ell$-primary torsion
and commutes with $\tau$. If $T=\tau\eta-\eta$, then
$T\ne0$ and
\[
                 0=(\tau^\ell-1)\eta=\ell U T,
                 \qquad \ell T=0.
\]
Since $\tau$ fixes $J[\ell]$, it fixes $T$. Its orbit on
$\eta$ is therefore the full coset
$\eta+\langle T\rangle$, of size $\ell$. Every point of
that coset lies on $\Theta$, because $\tau$ is a Frobenius
power, and also on $\Theta-T$.

The two irreducible divisors $\Theta$ and $\Theta-T$ are
distinct: a translation preserving a principal theta divisor
lies in the trivial kernel of its principal polarization.
Their intersection number is $\Theta^2=2$, so they cannot
contain $\ell\ge3$ distinct common points. This contradiction
proves $\sigma\eta=\eta$, which is (C2.27). No scalar or
semisimple Frobenius hypothesis was used.

For the curve in A.C2.11, direct counts give
$\#C(\F_4)=5$ and $\#C(\F_{16})=15$. The Frobenius
polynomial is $P(X)=X^4-X^2+16$. Exact reduction gives
\[
          X^{12}-1\equiv3(235X^2-1115)\pmod{P(X)},
\]
so $\pi^{12}$ fixes $J[3]$. By (C2.27) every three-primary
theta point is fixed by $\pi^{12}$. The further exact identity
\[
\begin{aligned}
 &(X^{12}-1)(176+47X^2)+27\cdot41440\\
 &\quad=P(X)\bigl(47X^{10}+223X^8-529X^6-4097X^4
                              +4367X^2+69919\bigr)
\end{aligned}                                                     \tag{C2.28}
\]
then kills it by $27\cdot41440$. Since $41440$ is prime to
three, its order divides $27$. These are polynomial
endomorphism identities; nontrivial Jordan reduction causes
no additional assumption. The counts and both identities are
verified in \path{code/A_theta_frobenius_certificate.py}.
Root independently checked the proof and integer remainders
in \path{audit/A_theta_prime_power_root_review.md}.
\end{proof}

Thus repeated division does not force theta points of arbitrarily
high prescribed prime-power order. For each fixed prime, A.C2.12
reduces that route to a finite set. Neither prime-to-two torsion
density nor surjectivity onto selected primary components forces
all other primary components to vanish. The required existence
statement for arbitrary genus-two curves remains open here.
